% GENERATED FILE. Edit canonical lessons, then run npm run generate:books.
% canonical-source-hash: fnv1a64:0780fff064a21883
% canonical-lessons: GU-C536-muhurt, GU-C536-coghdiyu, GU-C536-pancag, GU-C536-kundli, GU-C536-bhavi

\chapter{An Auspicious Hour}
\label{ch:gu-an-auspicious-hour}
\hlchaptermodality{536}

\begin{chapteropening}
\textbf{By the end of this chapter:} \emph{I can talk about choosing a lucky day and time.}

Complete the last lesson of chapter 536: I can talk about choosing a lucky day and time.
\end{chapteropening}

\section[\texorpdfstring{\gu{મુહૂર્ત} (muhūrt)}{મુહૂર્ત (muhūrt)}]{\gu{મુહૂર્ત} (muhūrt) — an auspicious moment}

\label{lesson:GU-C536-muhurt}



\begin{quote}
Before the new one: say the Gujarati for a planet, then the Gujarati for to make out.
\end{quote}

\subsection*{You'll want to know: \gu{મુહૂર્ત}}
\textbf{\gu{મુહૂર્ત}} — \emph{muhūrt} — \textquotedblleft{}an auspicious moment\textquotedblright{}.

\begin{quote}
\gu{જો} \gu{સારું} \gu{મુહૂર્ત} \gu{હશે} \gu{તો} \gu{અમે} \gu{રવિવારે} \gu{નવા} \gu{ઘરમાં} \gu{જઈશું}.
\end{quote}

\begin{quote}
\emph{jo sārũ muhūrt haśe to ame ravivāre navā gharmā̃ jaīśũ.}
\end{quote}

\begin{quote}
If there is a good hour, we will move into the new house on Sunday.
\end{quote}

A neuter noun.

\begin{cousinweb}
\textbf{\gu{મુહૂર્ત}} (\emph{muhūrt}) is a tatsama, taken from Sanskrit \emph{muhūrta}, \textquotedblleft{}a moment, a short space of time\textquotedblright{} (Monier-Williams); in Gujarati it is the right moment to begin.
\end{cousinweb}

\begin{grammarlens}[title={What you've built}]
One more word for saying what happens if.
\end{grammarlens}

\subsection*{Your turn}
\begin{itemize}
\raggedright
  \item \emph{Say it:} \emph{muhūrt}
  \item \emph{Say it:} \emph{muhūrt}, once more
  \item \emph{Say it:} \emph{jo sārũ muhūrt haśe to ame ravivāre navā gharmā̃ jaīśũ.} — If there is a good hour, we will move into the new house on Sunday.
  \item \emph{Recall:} say the Gujarati for a planet, then the Gujarati for to make out, then say \emph{muhūrt} again
\end{itemize}

\begin{tcolorbox}[breakable,colback=teal!4,colframe=teal!35!black,title={Before you move on}]
What does \gu{મુહૂર્ત} mean? (\textquotedblleft{}An auspicious moment\textquotedblright{}.) Say it once more.
\end{tcolorbox}
\section[\texorpdfstring{\gu{ચોઘડિયું} (coghḍiyũ)}{ચોઘડિયું (coghḍiyũ)}]{\gu{ચોઘડિયું} (coghḍiyũ) — a choghadiyu (a lucky or unlucky part of the day)}

\label{lesson:GU-C536-coghdiyu}



\begin{quote}
Before the new one: say the Gujarati for to make out, then the Gujarati for an auspicious moment.
\end{quote}

\subsection*{You'll want to know: \gu{ચોઘડિયું}}
\textbf{\gu{ચોઘડિયું}} — \emph{coghḍiyũ} — \textquotedblleft{}a choghadiyu (a lucky or unlucky part of the day)\textquotedblright{}.

\begin{quote}
\gu{જો} \gu{ચોઘડિયું} \gu{સારું} \gu{હશે} \gu{તો} \gu{દાદા} \gu{સવારે} \gu{દુકાન} \gu{ખોલશે}.
\end{quote}

\begin{quote}
\emph{jo coghḍiyũ sārũ haśe to dādā savāre dukān kholśe.}
\end{quote}

\begin{quote}
If the time of day is lucky, my grandfather will open the shop in the morning.
\end{quote}

A neuter noun.

\begin{cousinweb}
\textbf{\gu{ચોઘડિયું}} (\emph{coghḍiyũ}) is usually explained as built from \emph{co-}, \textquotedblleft{}four\textquotedblright{}, and \emph{ghaḍī}, an old unit of time, from Sanskrit \emph{ghaṭikā} (Monier-Williams): the day is cut into such periods, each good or bad for beginning a task.
\end{cousinweb}

\begin{grammarlens}[title={What you've built}]
One more word for saying what happens if.
\end{grammarlens}

\subsection*{Your turn}
\begin{itemize}
\raggedright
  \item \emph{Say it:} \emph{coghḍiyũ}
  \item \emph{Say it:} \emph{coghḍiyũ}, once more
  \item \emph{Say it:} \emph{jo coghḍiyũ sārũ haśe to dādā savāre dukān kholśe.} — If the time of day is lucky, my grandfather will open the shop in the morning.
  \item \emph{Recall:} say the Gujarati for to make out, then the Gujarati for an auspicious moment, then say \emph{coghḍiyũ} again
\end{itemize}

\begin{tcolorbox}[breakable,colback=teal!4,colframe=teal!35!black,title={Before you move on}]
What does \gu{ચોઘડિયું} mean? (\textquotedblleft{}A choghadiyu (a lucky or unlucky part of the day)\textquotedblright{}.) Say it once more.
\end{tcolorbox}
\section[\texorpdfstring{\gu{પંચાંગ} (pancā̃g)}{પંચાંગ (pancā̃g)}]{\gu{પંચાંગ} (pancā̃g) — a Hindu almanac}

\label{lesson:GU-C536-pancag}



\begin{quote}
Before the new one: say the Gujarati for an auspicious moment, then the Gujarati for a choghadiyu.
\end{quote}

\subsection*{You'll want to know: \gu{પંચાંગ}}
\textbf{\gu{પંચાંગ}} — \emph{pancā̃g} — \textquotedblleft{}a Hindu almanac\textquotedblright{}.

\begin{quote}
\gu{જો} \gu{તમે} \gu{પંચાંગ} \gu{જોશો} \gu{તો} \gu{તમને} \gu{તહેવારની} \gu{તારીખ} \gu{મળશે}.
\end{quote}

\begin{quote}
\emph{jo tame pancā̃g jośo to tamne tahevārnī tārīkh maḷśe.}
\end{quote}

\begin{quote}
If you look in the almanac, you'll find the date of the festival.
\end{quote}

A neuter noun.

\begin{cousinweb}
\textbf{\gu{પંચાંગ}} (\emph{pancā̃g}) is a tatsama, taken from Sanskrit \emph{pañcāṅga}, \textquotedblleft{}having five parts\textquotedblright{} (Monier-Williams); the almanac gives five things for each day.
\end{cousinweb}

\begin{grammarlens}[title={What you've built}]
One more word for saying what happens if.
\end{grammarlens}

\subsection*{Your turn}
\begin{itemize}
\raggedright
  \item \emph{Say it:} \emph{pancā̃g}
  \item \emph{Say it:} \emph{pancā̃g}, once more
  \item \emph{Say it:} \emph{jo tame pancā̃g jośo to tamne tahevārnī tārīkh maḷśe.} — If you look in the almanac, you'll find the date of the festival.
  \item \emph{Recall:} say the Gujarati for an auspicious moment, then the Gujarati for a choghadiyu, then say \emph{pancā̃g} again
\end{itemize}

\begin{tcolorbox}[breakable,colback=teal!4,colframe=teal!35!black,title={Before you move on}]
What does \gu{પંચાંગ} mean? (\textquotedblleft{}A Hindu almanac\textquotedblright{}.) Say it once more.
\end{tcolorbox}
\section[\texorpdfstring{\gu{કુંડળી} (kunḍḷī)}{કુંડળી (kunḍḷī)}]{\gu{કુંડળી} (kunḍḷī) — a birth chart}

\label{lesson:GU-C536-kundli}



\begin{quote}
Before the new one: say the Gujarati for a choghadiyu, then the Gujarati for a Hindu almanac.
\end{quote}

\subsection*{You'll want to know: \gu{કુંડળી}}
\textbf{\gu{કુંડળી}} — \emph{kunḍḷī} — \textquotedblleft{}a birth chart\textquotedblright{}.

\begin{quote}
\gu{જો} \gu{બંનેની} \gu{કુંડળી} \gu{મળશે} \gu{તો} \gu{લગ્ન} \gu{નક્કી} \gu{થશે}.
\end{quote}

\begin{quote}
\emph{jo bannenī kunḍḷī maḷśe to lagna nakkī thaśe.}
\end{quote}

\begin{quote}
If both their birth charts match, the wedding will be fixed.
\end{quote}

A feminine noun.

\begin{cousinweb}
\textbf{\gu{કુંડળી}} (\emph{kunḍḷī}) goes back to Sanskrit \emph{kuṇḍala}, \textquotedblleft{}a ring, a circle\textquotedblright{} (Monier-Williams); the chart of the heavens at one's birth is drawn as a figure.
\end{cousinweb}

\begin{grammarlens}[title={What you've built}]
One more word for saying what happens if.
\end{grammarlens}

\subsection*{Your turn}
\begin{itemize}
\raggedright
  \item \emph{Say it:} \emph{kunḍḷī}
  \item \emph{Say it:} \emph{kunḍḷī}, once more
  \item \emph{Say it:} \emph{jo bannenī kunḍḷī maḷśe to lagna nakkī thaśe.} — If both their birth charts match, the wedding will be fixed.
  \item \emph{Recall:} say the Gujarati for a choghadiyu, then the Gujarati for a Hindu almanac, then say \emph{kunḍḷī} again
\end{itemize}

\begin{tcolorbox}[breakable,colback=teal!4,colframe=teal!35!black,title={Before you move on}]
What does \gu{કુંડળી} mean? (\textquotedblleft{}A birth chart\textquotedblright{}.) Say it once more.
\end{tcolorbox}
\section[\texorpdfstring{\gu{ભાવિ} (bhāvi)}{ભાવિ (bhāvi)}]{\gu{ભાવિ} (bhāvi) — the future}

\label{lesson:GU-C536-bhavi}



\begin{quote}
Before the new one: say the Gujarati for a Hindu almanac, then the Gujarati for a birth chart.
\end{quote}

\subsection*{You'll want to know: \gu{ભાવિ}}
\textbf{\gu{ભાવિ}} — \emph{bhāvi} — \textquotedblleft{}the future\textquotedblright{}.

\begin{quote}
\gu{જો} \gu{તમે} \gu{મહેનત} \gu{કરશો} \gu{તો} \gu{તમારું} \gu{ભાવિ} \gu{સારું} \gu{બનશે}.
\end{quote}

\begin{quote}
\emph{jo tame mahenat karśo to tamārũ bhāvi sārũ banśe.}
\end{quote}

\begin{quote}
If you work hard, your future will be good.
\end{quote}

A neuter noun.

\begin{cousinweb}
\textbf{\gu{ભાવિ}} (\emph{bhāvi}) is a tatsama, from Sanskrit \emph{bhāvin}, \textquotedblleft{}about to be, future\textquotedblright{} (Monier-Williams).
\end{cousinweb}

\begin{grammarlens}[title={What you've built}]
One more word for saying what happens if.
\end{grammarlens}

\subsection*{Your turn}
\begin{itemize}
\raggedright
  \item \emph{Say it:} \emph{bhāvi}
  \item \emph{Say it:} \emph{bhāvi}, once more
  \item \emph{Say it:} \emph{jo tame mahenat karśo to tamārũ bhāvi sārũ banśe.} — If you work hard, your future will be good.
  \item \emph{Recall:} say the Gujarati for a Hindu almanac, then the Gujarati for a birth chart, then say \emph{bhāvi} again
\end{itemize}

\begin{tcolorbox}[breakable,colback=teal!4,colframe=teal!35!black,title={Before you move on}]
What does \gu{ભાવિ} mean? (\textquotedblleft{}The future\textquotedblright{}.) Say it once more.
\end{tcolorbox}
